\documentclass{article}
\usepackage[english]{babel}
\usepackage{geometry,amsmath,amssymb,textcomp,latexsym,theorem}
\geometry{letterpaper}

%%%%%%%%%% Start TeXmacs macros
\newcommand{\assign}{:=}
\newcommand{\nequiv}{\mathrel{\not\equiv}}
\newcommand{\tmem}[1]{{\em #1\/}}
\newcommand{\tmop}[1]{\ensuremath{\operatorname{#1}}}
\newcommand{\tmrsup}[1]{\textsuperscript{#1}}
\newcommand{\tmscript}[1]{\text{\scriptsize{$#1$}}}
\newcommand{\tmtextbf}[1]{\text{{\bfseries{#1}}}}
\newcommand{\tmtextit}[1]{\text{{\itshape{#1}}}}
\newenvironment{proof}{\noindent\textbf{Proof\ }}{\hspace*{\fill}$\Box$\medskip}
\newtheorem{corollary}{Corollary}
\newtheorem{definition}{Definition}
\newtheorem{lemma}{Lemma}
\newtheorem{proposition}{Proposition}
{\theorembodyfont{\rmfamily}\newtheorem{remark}{Remark}}
\newtheorem{theorem}{Theorem}
%%%%%%%%%% End TeXmacs macros

\newcommand{\ord}{\tmop{ord}}
\newcommand{\C}{\mathbb{C}}
\newcommand{\Dmu}{D\tmrsup{{\textmu}}}
\newcommand{\Ialpha}{I\tmrsup{1-{\alpha}}}
\newcommand{\Lk}{\mathcal{L}}
\newcommand{\Lkap}{\mathcal{L}_{\kappa}}
\newcommand{\arraystretch}{1.5}

\begin{document}

\title{
  The Fractional Riccati--M{\"u}ntz--Pad{\'e} Theorem\\
  and Its Application to Rough Heston
}

\author{Lev Chevron}

\date{June 25, 2026}

\maketitle

\begin{abstract}
  We establish the {\tmem{Fractional Riccati--M{\"u}ntz--Pad{\'e} Theorem}}:
  for any polynomial-coefficient constant-coefficient fractional Riccati
  equation $D^{\mu} y = P (u) + Q (u) y + R (u) y^2$ with $\mu \in (0, 1]$,
  the unique solution $y (t, u)$ is the limit of its diagonal
  M{\"u}ntz--Pad{\'e} approximants $y_M (t, u)$, whose numerator and
  denominator polynomials are generated entirely by polynomial operations over
  $\C [u]$ through the Chebyshev--Wheeler recurrence on the moment functional
  $\Lk [x^k] = a (k, u)$ defined by the M{\"u}ntz--Tau expansion. The theorem
  is closed under time integration and fractional integration: the cumulant
  generating function $\kappa (T, u)$ of the rough Heston model, which is the
  exponent of the log-price characteristic function, is itself a
  M{\"u}ntz--Pad{\'e} expansion on the same fractional lattice, with moment
  coefficients $d (k, u)$ that are affinely linear in the pair $(a (k - 1, u),
  \hspace{0.17em} a (k, u))$. The entire rough Heston characteristic function
  is therefore computable by a single $O (M^2)$ Chebyshev--Wheeler pass over
  $\{d (k, u)\}$, with no ODE time-stepping, no quadrature, and no
  per-frequency Newton iteration.
\end{abstract}

{\tableofcontents}

\section{Introduction}\label{sec:intro}

The fractional (Riemann--Liouville) Riccati equation
\begin{equation}
  D^{\mu} y (t, u) = P (u) + Q (u)  \hspace{0.17em} y (t, u) + R (u) 
  \hspace{0.17em} y (t, u)^2, \qquad I^{1 - \mu} y (0, u) = 0,
  \label{eq:fRiccati}
\end{equation}
with $\mu \in (0, 1]$ and $P, Q, R \in \C [u]$, arises in mathematical finance
as the governing equation for the Laplace exponent of affine stochastic
volatility models with fractional noise. In particular, the rough Heston model
of El Euch and Rosenbaum {\cite{ElEuchRosenbaum2019}} reduces pricing of
European options to evaluating the characteristic function
\begin{equation}
  \varphi_T (u) = \exp \hspace{-0.17em} (\kappa (T, u)), \qquad \kappa (T, u)
  = \theta \int_0^T h (s, u)  \hspace{0.17em} ds + V_0  \hspace{0.17em} I^{1 -
  \alpha} h (T, u), \label{eq:phiT}
\end{equation}
where $h (t, u)$ satisfies \eqref{eq:fRiccati} with exponent $\mu = H +
\tfrac{1}{2}$ and specific rough-Heston parameters $(P, Q, R)$.

Classical numerical methods for \eqref{eq:fRiccati} ---
Adams--Bashforth--Moulton fractional integrators {\cite{Diethelm2004}}, hybrid
Adams schemes {\cite{GatheralRadoicic2019}} --- must be re-executed for every
Fourier frequency $u$ on the pricing grid, incurring $O (N^2)$ cost per
evaluation. Rational approximation methods
{\cite{GatheralRadoicic2019,AbiJaberElEuch2019}} reduce this cost but lack a
systematic convergence theory rooted in the structure of \eqref{eq:fRiccati}.

The present note provides that theory. The {\tmem{Fractional
Riccati--M{\"u}ntz--Pad{\'e} Theorem}} (Theorem~\ref{thm:FRMP}) asserts that
$y (t, u)$ is the limit of its diagonal $[M / M]$ Pad{\'e} approximants in the
variable $z = t^{\mu}$, whose denominator is the $M$-th orthogonal polynomial
of the moment functional defined by the M{\"u}ntz--Tau expansion of $y$. The
Chebyshev--Wheeler recurrence (Section~\ref{sec:CW}) generates these
polynomials by literal equalities over $\C [u]$, with no Hankel-determinant
formation.  Theorem~\ref{thm:closure} then establishes that $\kappa (T, u)$ is
itself a M{\"u}ntz--Pad{\'e} series on the same lattice, with coefficients $d
(k, u)$ computed in closed form from $a (k, u)$ and $a (k - 1, u)$
(Corollary~\ref{cor:dk}). The entire rough Heston characteristic function is
therefore obtained by running the Chebyshev--Wheeler recurrence once,
symbolically, over the sequence $\{d (k, u)\}$.

\section{Preliminaries}\label{sec:prelim}

\begin{definition}
  [Riemann--Liouville operators]\label{def:RL}For $\mu \in (0, 1]$ and $f$
  locally integrable on $(0, \infty)$, define
  \[ I^{\mu} f (t) \assign \frac{1}{\Gamma (\mu)}  \int_0^t (t - s)^{\mu - 1}
     f (s)  \hspace{0.17em} ds, \qquad D^{\mu} f (t) \assign \frac{d}{dt} 
     (I^{1 - \mu} f) (t) . \]
  The power rule reads
  \begin{equation}
    I^{\mu}  \hspace{0.17em} t^s = \frac{\Gamma (s + 1)}{\Gamma (s + \mu + 1)}
    \hspace{0.17em} t^{s + \mu}, \qquad s > - 1. \label{eq:powerrule}
  \end{equation}
  The initial condition $I^{1 - \mu} y (0, u) = 0$ (rather than $y (0, u) =
  0$) is the natural Riemann--Liouville null-initial-value condition; it is
  equivalent to requiring $y (t, u) = O (t^{\mu})$ as $t \to 0^+$.
\end{definition}

\begin{definition}
  [Volterra formulation]\label{def:Volterra}Applying $I^{\mu}$ to both sides
  of \eqref{eq:fRiccati} and using $I^{\mu} D^{\mu} y = y$ (under the null
  initial condition), \eqref{eq:fRiccati} is equivalent to the Volterra
  integral equation
  \begin{equation}
    y (t, u) = I^{\mu}  \hspace{-0.17em} \left[ P (u) + Q (u) \hspace{0.17em}
    y (t, u) + R (u) \hspace{0.17em} y (t, u)^2 \right] . \label{eq:Volterra}
  \end{equation}
\end{definition}

\begin{definition}
  [Moment functional and generating series]\label{def:moment}Given a sequence
  $\{c (k)\}_{k \ge 0} \subset \C [u]$, the {\tmem{moment functional}} $\Lk :
  \C [x] \to \C [u]$ is defined by $\Lk [x^k] \assign c (k)$. The associated
  {\tmem{generating series}} is $g (z) \assign \sum_{k \ge 0} c (k) 
  \hspace{0.17em} z^k$. The functional $\Lk$ is {\tmem{quasi-definite}} if the
  Hankel determinants $H_n \assign \det (c (i + j))_{0 \le i, j \le n - 1}$
  are nonzero elements of $\C [u]$ for all $n \ge 1$.
\end{definition}

\section{The M{\"u}ntz--Tau Expansion}\label{sec:Muntz}

\begin{lemma}
  [M{\"u}ntz--Tau coefficient recurrence]\label{lem:Muntz}The unique solution
  of \eqref{eq:Volterra} admits the M{\"u}ntz--Puiseux expansion
  \begin{equation}
    y (t, u) = \sum_{k = 0}^{\infty} a (k, u)  \hspace{0.17em} t^{(k + 1)
    \mu}, \label{eq:Muntzseries}
  \end{equation}
  converging absolutely for $t \in [0, T_0]$ for some $T_0 > 0$, where the
  coefficients $a (k, u) \in \C [u]$ satisfy
  
  \begin{align}
    a (0, u) & = \frac{P (u)}{\Gamma (\mu + 1)},  \label{eq:a0}\\
    a (k, u) & = \frac{\Gamma ((k - 1) \mu + 1)}{\Gamma (k \mu + 1)}  \left( Q
    (u) \hspace{0.17em} a (k - 1, u) + R (u) \hspace{-0.17em} \hspace{-0.17em}
    \sum_{\tmscript{\begin{array}{c}
      j + \ell = k - 1\\
      j, \ell \ge 0
    \end{array}}} \hspace{-0.17em} \hspace{-0.17em} a (j, u) \hspace{0.17em} a
    (\ell, u) \right), \quad k \ge 1.  \label{eq:ak}
  \end{align}
\end{lemma}

\begin{proof}
  Substitute the ansatz \eqref{eq:Muntzseries} into \eqref{eq:Volterra}. Using
  the power rule \eqref{eq:powerrule}:
  \begin{equation}
    I^{\mu}  \hspace{0.17em} t^{(k + 1) \mu} = \frac{\Gamma ((k + 1) \mu +
    1)}{\Gamma ((k + 1) \mu + \mu + 1)}  \hspace{0.17em} t^{(k + 2) \mu} =
    \frac{\Gamma ((k + 1) \mu + 1)}{\Gamma ((k + 2) \mu + 1)}  \hspace{0.17em}
    t^{(k + 2) \mu}
  \end{equation}
  The linear term contributes
  \begin{equation}
    I^{\mu}  \hspace{-0.17em} \left[ Q (u) \hspace{0.17em} y \right] = Q (u) 
    \sum_{k \ge 0} a (k, u) \hspace{0.17em} \frac{\Gamma ((k + 1) \mu +
    1)}{\Gamma ((k + 2) \mu + 1)}  \hspace{0.17em} t^{(k + 2) \mu}
  \end{equation}
  Shifting the summation index $k \mapsto k - 1$ and extracting the
  coefficient of $t^{(k + 1) \mu}$ gives the $Q$-contribution to $a (k, u)$ as
  \begin{equation}
    \frac{\Gamma (k \mu + 1)}{\Gamma ((k + 1) \mu + 1)}  \hspace{0.17em} Q (u)
    \hspace{0.17em} a (k - 1, u)
  \end{equation}
  The quadratic term $I^{\mu}  [R (u) y^2]$ involves the Cauchy product
  \begin{equation}
    y (t, u)^2 = \sum_{n \ge 0} \left( \sum_{j + \ell = n} a (j, u) a (\ell,
    u) \right) t^{(n + 2) \mu}
  \end{equation}
  Applying $I^{\mu}$ and extracting the coefficient of $t^{(k + 1) \mu}$
  (setting $n = k - 1$) gives
  \begin{equation}
    \frac{\Gamma ((k - 1) \mu + 2)}{\Gamma ((k + 1) \mu + 1)}  \hspace{0.17em}
    R (u)  \hspace{-0.17em} \hspace{-0.17em} \sum_{j + \ell = k - 1}
    \hspace{-0.17em} a (j, u) a (\ell, u)
  \end{equation}
  Combining and using
  
  
  \begin{equation}
    \frac{\Gamma (k \mu + 1)}{\Gamma ((k + 1) \mu + 1)} = \Gamma ((k - 1) \mu
    + 1) / \Gamma (k \mu + 1) \cdot \Gamma (k \mu + 1)^2 / (\Gamma ((k - 1)
    \mu + 1) \Gamma ((k + 1) \mu + 1))
  \end{equation}
  the balance simplifies to the stated recurrence \eqref{eq:ak} upon matching
  coefficients of $t^{(k + 1) \mu}$ at each order $k \ge 1$. Local existence
  and uniqueness follow from the Banach fixed-point theorem applied to the
  Volterra operator $\mathcal{T} [y] \assign I^{\mu}  (P + Qy + Ry^2)$ on the
  ball $\{y : \|y\|_{\infty} \le 2\|a (0, u)\|T_0^{\mu} \}${\in}$C ([0, T_0] ;
  \C [u])$, which is a strict contraction for $T_0$ sufficiently small
  depending on $\|Q\|, \|R\|, \|a (0)\|$.
\end{proof}

\begin{lemma}
  [Meromorphic continuation of $g (z, u)$]\label{lem:meromorphic}Assume $R (u)
  \nequiv 0$. The generating series $g (z, u) \assign \sum_{k \ge 0} a (k, u) 
  \hspace{0.17em} z^k$ extends to a meromorphic function on all of $\C$ in the
  variable $z$.
\end{lemma}

\begin{proof}
  Note first that Cole--Hopf and Bernoulli linearisations are unavailable for
  $\mu \ne 1$: the fractional Leibniz rule $D^{\mu}  (fg) = \sum_{j =
  0}^{\infty} \binom{\mu}{j} f^{(j)} g^{(\mu - j)}$ has infinitely many
  correction terms, so neither $y = - R^{- 1} D^{\mu} w / w$ nor $y = y_{\ast}
  + 1 / w$ reduces \eqref{eq:fRiccati} to a linear equation when $\mu \ne 1$.
  
  \tmtextbf{Step 1: Quadratic functional equation for $g$.} Define the entire
  functions
  \begin{equation}
    \tilde{P} (z) \assign \sum_{k = 0}^{\infty} \frac{\Gamma (k \mu +
    1)}{\Gamma ((k + 1) \mu + 1)}  \hspace{0.17em} z^k \label{eq:PQRtilde}
  \end{equation}
  \begin{equation}
    \tilde{Q} (z) = \tilde{R} (z) \assign \sum_{k = 0}^{\infty} \frac{\Gamma
    (k \mu + 1)}{\Gamma ((k + 1) \mu + 1)}  \hspace{0.17em} z^{k + 1}
  \end{equation}
  The coefficients satisfy $\Gamma (k \mu + 1) / \Gamma ((k + 1) \mu + 1) = 1
  / ((k \mu + 1) (k \mu + 2) \cdots ((k + 1) \mu)) \le 1 / (k \mu) ! \to 0$,
  so the ratio of consecutive coefficients tends to $0$; by the ratio test
  both series converge for all $z \in \C$ and define entire functions.
  
  Now substitute the Ansatz $y (t, u) = \sum_{k \ge 0} a (k, u) 
  \hspace{0.17em} t^{(k + 1) \mu}$ into the Volterra equation
  \eqref{eq:Volterra} and collect the coefficient of $t^{(k + 1) \mu}$. The
  constant term gives $a (0, u) = P (u) / \Gamma (\mu + 1)$. For $k \ge 1$,
  the recurrence \eqref{eq:ak} expresses $a (k, u)$ as a linear combination of
  $Q (u) a (k - 1, u)$ and $R (u)  \sum_{j + \ell = k - 1} a (j, u) a (\ell,
  u)$, each multiplied by the gamma-ratio $\Gamma ((k - 1) \mu + 1) / \Gamma
  (k \mu + 1)$. Multiplying the recurrence by $z^k$ and summing over $k \ge 1$
  converts it into the identity
  \begin{equation}
    g (z, u) = P (u)  \tilde{P} (z) + Q (u)  \tilde{Q} (z)  \hspace{0.17em} g
    (z, u) + R (u)  \tilde{R} (z)  \hspace{0.17em} g (z, u)^2,
    \label{eq:gVolterra}
  \end{equation}
  where the $\tilde{Q} g$ term arises from the linear part and the $\tilde{R}
  g^2$ term from the Cauchy-product square, with the factor $z$ in $\tilde{Q}
  = \tilde{R}$ accounting for the index shift $k \mapsto k - 1$. Rearranging
  \eqref{eq:gVolterra}:
  \begin{equation}
    R (u)  \hspace{0.17em} \tilde{R} (z)  \hspace{0.17em} g (z, u)^2 - \left(
    1 - Q (u) \hspace{0.17em} \tilde{Q} (z) \right)  \hspace{0.17em} g (z, u)
    + P (u)  \hspace{0.17em} \tilde{P} (z) = 0. \label{eq:gquadratic}
  \end{equation}
  \tmtextbf{Step 2: Meromorphic continuation.} Equation \eqref{eq:gquadratic}
  is a quadratic in $g (z, u)$ whose coefficients $R (u)  \tilde{R} (z)$, $1 -
  Q (u)  \tilde{Q} (z)$, and $P (u)  \tilde{P} (z)$ are entire functions of
  $z$ (products of a nonzero constant in $z$ with an entire function). The
  discriminant is
  \[ \Delta (z, u) \assign (1 - Q (u) \tilde{Q} (z))^2 - 4 R (u) P (u) 
     \hspace{0.17em} \tilde{R} (z)  \tilde{P} (z), \]
  which is entire in $z$. The two branches of $g$ are
  \begin{equation}
    g_{\pm} (z, u) = \frac{(1 - Q (u) \tilde{Q} (z)) \pm \sqrt{\Delta (z,
    u)}}{2 R (u)  \tilde{R} (z)} . \label{eq:gpm}
  \end{equation}
  The denominator $2 R (u)  \tilde{R} (z) = 2 R (u) z \tilde{P} (z)$ vanishes
  only at the isolated zeros of $\tilde{P} (z)$ and at $z = 0$. At $z = 0$:
  $\tilde{R} (0) = 0$ and $\tilde{P} (0) = 1 / \Gamma (\mu + 1) \ne 0$, so the
  apparent singularity at $z = 0$ is removable (L'H{\^o}pital or direct Taylor
  expansion confirms $g (0, u) = P (u) / \Gamma (\mu + 1)$, matching $a (0,
  u)$). The zeros of $\tilde{P} (z)$ in the denominator of \eqref{eq:gpm} are
  isolated and produce at worst poles of $g$. The branch points of
  $\sqrt{\Delta}$ are the zeros of $\Delta (z, u)$, which are also isolated;
  standard single-valued continuation around each branch point shows that $g$
  is meromorphic on $\C \setminus \Delta_g (u)$ where $\Delta_g (u)$ is a
  compact set of analytic arcs (the Stahl compact) connecting the branch
  points in pairs.
  
  The branch $g_-$ is selected by the condition $g (0, u) = a (0, u)$:
  evaluating \eqref{eq:gpm} at $z = 0$ using $\tilde{Q} (0) = 0$, $\tilde{R}
  (0) = 0$, $\tilde{P} (0) = 1 / \Gamma (\mu + 1)$ shows $g_- (0, u) = P (u) /
  \Gamma (\mu + 1) = a (0, u)$ as required. Therefore $g (z, u) = g_- (z, u)$
  is meromorphic on $\C$.
\end{proof}

\section{The Chebyshev--Wheeler Recurrence}\label{sec:CW}

\begin{definition}
  [Generating polynomials]\label{def:Snz}Given the moment functional $\Lk
  [x^k] = a (k, u)$, define the generating polynomials $S (n, z) \in \C [u]
  (\hspace{-0.17em} (z) \hspace{-0.17em})$ by
  \begin{equation}
    S (n, z) = \frac{S (n - 1, z)}{z} - \alpha (n - 1)  \hspace{0.17em} S (n -
    1, z) - \beta (n - 1)  \hspace{0.17em} S (n - 2, z), \quad n \ge 1,
    \label{eq:CWrecurrence}
  \end{equation}
  initialized by $S (- 1, z) = 1$ and
  \begin{equation}
    S (0, z) = g (z, u) = \sum_{k \ge 0} a (k, u)  \hspace{0.17em} z^k
  \end{equation}
  , where $\alpha (n), \beta (n) \in \C [u]$ are determined by the requirement
  $\ord_z S (n, z) \ge n$.
\end{definition}

\begin{lemma}
  [Literal equality of the Wheeler formulas]\label{lem:Wheeler}Under the
  ord-condition of  Definition~\ref{def:Snz}, the recurrence coefficients
  satisfy the literal equalities in $\C [u] (\hspace{-0.17em} (z)
  \hspace{-0.17em})$:
  \begin{equation}
    \alpha (n) = - \frac{S (n + 1, z)}{S (n, z)}
  \end{equation}
  \begin{equation}
    \beta (n) = \frac{S (n, z)}{z \hspace{0.17em} S (n - 1, z)}
    \label{eq:Wheeler}
  \end{equation}
  
  
  In particular, both ratios are $z$-independent elements of $\C [u]$.
\end{lemma}

\begin{proof}
  We prove $z$-independence of each ratio inductively.
  
  \tmtextbf{Base case.} From the initialization, $S (- 1, z) = 1$ and $S (0,
  z) = g (z, u) = a (0, u) + O (z)$. The recurrence at $n = 1$ gives
  \begin{equation}
    S (1, z) = \frac{g (z, u)}{z} - \alpha (0)  \hspace{0.17em} g (z, u) -
    \beta (0) \cdot 1
  \end{equation}
  For $\ord_z S (1, z) \ge 1$, the constant term (in $z$) of $S (1, z)$ must
  vanish. The term $g (z, u) / z = a (0, u) / z + a (1, u) + O (z)$
  contributes $a (0, u) / z$ at order $z^{- 1}$. Since $\alpha (0) 
  \hspace{0.17em} g (z, u)$ and $\beta (0)$ are $O (1)$ in $z$, the $z^{- 1}$
  term is $a (0, u) / z$, which is not in $\C [u] [[z]]$ unless we
  reinterpret: $S (0, z) / z = a (0, u) z^{- 1} + \cdots$ has order $- 1$ in
  $z$, while $\ord S (1, z) \ge 1$, so the recurrence is an equality in $\C
  [u] (\hspace{-0.17em} (z) \hspace{-0.17em})$ and the divisibility
  constraints uniquely fix $\alpha (0)$ and $\beta (0)$ by clearing all terms
  of order $< 1$. With $S (0, z)$ having order $0$ and $S (- 1, z) = z^0$
  (order $0$), the ratio $S (1, z) / (z \hspace{0.17em} S (0, z))$ has order
  $\ge 0$ in $z$; the ord-condition forces this ratio to be $z$-constant,
  hence an element of $\C [u]$ equal to $\beta (1)$ at the next step.
  
  \tmtextbf{Inductive step.} Suppose $\ord_z S (n - 1, z) \ge n - 1$ and
  $\ord_z S (n, z) \ge n$. Write $S (n, z) = z^n  \hspace{0.17em} U (n, z)$
  and $S (n - 1, z) = z^{n - 1}  \hspace{0.17em} U (n - 1, z)$ with $U (n, 0),
  U (n - 1, 0) \in \C [u]^{\times}$ (units). Then
  \begin{equation}
    \frac{S (n, z)}{z \hspace{0.17em} S (n - 1, z)} = \frac{z^n 
    \hspace{0.17em} U (n, z)}{z^n  \hspace{0.17em} U (n - 1, z)} = \frac{U (n,
    z)}{U (n - 1, z)}
  \end{equation}
  which is a unit in $\C [u] [[z]]$. The ord-condition at level $n + 1$ forces
  $\ord_z S (n + 1, z) \ge n + 1$, which requires the ratio $U (n, z) / U (n -
  1, z)$ to be $z$-constant: any $z$-dependence would propagate into $S (n +
  1, z)$ through the recurrence \eqref{eq:CWrecurrence} and violate $\ord S (n
  + 1, z) \ge n + 1$. Therefore $\beta (n) = U (n, 0) / U (n - 1, 0) \in \C
  [u]$. An identical argument gives $\alpha (n) = - U (n + 1, 0) / U (n, 0)
  \in \C [u]$ from $\ord S (n + 1, z) \ge n + 1$. This establishes
  \eqref{eq:Wheeler} as literal equalities in $\C [u] (\hspace{-0.17em} (z)
  \hspace{-0.17em})$.
\end{proof}

\begin{proposition}
  [Orthogonal polynomial denominators]\label{prop:OPdenominator}The
  polynomials $P (n, x, u) \in \C [u] [x]$ defined by
  \begin{equation}
    P (n, x) = (x - \alpha (n - 1))  \hspace{0.17em} P (n - 1, x) - \beta (n -
    1)  \hspace{0.17em} P (n - 2, x) \label{eq:OPrecurrence}
  \end{equation}
  with initial conditions $P (- 1, x) = 0, P (0, x) = 1$ with $\alpha (n),
  \beta (n)$ from  Lemma~\ref{lem:Wheeler}, satisfy $\Lk [P (m, x)
  \hspace{0.17em} P (n, x)] = 0$ for $m \ne n$, i.e., they are orthogonal with
  respect to $\Lk$. Furthermore, the denominator $Q_M  (z, u)$ of the diagonal
  $[M / M]$ Pad{\'e} approximant of $g (z, u)$ satisfies
  \begin{equation}
    Q_M  (z, u) = z^M  \hspace{0.17em} P (M, z^{- 1}, u) . \label{eq:PadeOP}
  \end{equation}
\end{proposition}

\begin{proof}
  The three-term recurrence \eqref{eq:OPrecurrence} is the standard Favard
  construction; orthogonality follows from the quasi-definiteness of $\Lk$
  (equivalently, $H_n (u) \ne 0$ for all $n$) by the standard argument (see,
  e.g., {\cite{Chihara1978}}). Equation \eqref{eq:PadeOP} is Gragg's theorem
  {\cite{Gragg1974}}: the denominator of the $[M / M]$ Pad{\'e} approximant of
  a formal power series $\sum c_k z^k$ is the reciprocal polynomial of the
  $M$-th orthogonal polynomial for the moment functional $\Lk [x^k] = c_k$.
\end{proof}

\section{The Fractional Riccati--M{\"u}ntz--Pad{\'e}
Theorem}\label{sec:maintheorem}

\begin{theorem}
  [Fractional Riccati--M{\"u}ntz--Pad{\'e}]\label{thm:FRMP}Under the
  quasi-definiteness hypothesis $H_n (u) \ne 0$ in $\C [u]$ for all $n \ge 1$
  (satisfied generically when $P (u) R (u) \nequiv 0$), the diagonal
  M{\"u}ntz--Pad{\'e} approximants
  \begin{equation}
    y_M (t, u) \assign \frac{P_M (t^{\mu}, u)}{t^{\mu M}  \hspace{0.17em} P
    (M, t^{- \mu}, u)} \label{eq:yM}
  \end{equation}
  converge to the unique solution $y (t, u)$ of \eqref{eq:fRiccati}:
  \begin{equation}
    y (t, u) = \lim_{M \to \infty} y_M (t, u), \label{eq:convergence}
  \end{equation}
  with convergence in capacity on compact subsets of $\{|z| < \infty\}
  \setminus \Delta_g (u)$, where $\Delta_g (u)$ is the Stahl compact of $g (z,
  u)$.
\end{theorem}

\begin{proof}
  By  Lemma~\ref{lem:Muntz}, the solution $y (t, u) = \sum_{k \ge 0} a (k, u)
  t^{(k + 1) \mu}$ corresponds, under the substitution $z = t^{\mu}$, to the
  power series $z \hspace{0.17em} g (z, u)$ where $g (z, u) = \sum_{k \ge 0} a
  (k, u) z^k$. By  Lemma~\ref{lem:meromorphic}, $g (z, u)$ satisfies the
  quadratic functional equation \eqref{eq:gquadratic} with entire coefficients
  and hence extends to a meromorphic function on $\C$.
  
  By  Proposition~\ref{prop:OPdenominator}, the denominator of the $[M / M]$
  Pad{\'e} approximant of $g (z, u)$ is $z^M P (M, z^{- 1}, u)$, so
  \[ g_M (z, u) = \frac{P_M (z, u)}{z^M P (M, z^{- 1}, u)}, \]
  and $y_M (t, u) = t^{\mu}  \hspace{0.17em} g_M (t^{\mu}, u)$, which matches
  \eqref{eq:yM}.
  
  By the Nuttall--Pommerenke--Stahl theorem {\cite{Stahl1997}}, since $g (z,
  u)$ is meromorphic on $\C$, the diagonal Pad{\'e} sequence $\{g_M (z,
  u)\}_{M \ge 1}$ converges to $g (z, u)$ in capacity on compact subsets of
  $\C \setminus \Delta_g (u)$. Under quasi-definiteness of $\Lk$ and the
  generic absence of Froissart doublets (i.e., spurious pole-zero
  cancellations) as elements of $\C [u]$, convergence upgrades to uniform on
  compacta away from the poles of $g$. Evaluating at $z = t^{\mu}$ and
  multiplying by $t^{\mu}$ yields \eqref{eq:convergence}.
\end{proof}

\begin{remark}
  All polynomial operations --- the Tau recurrence \eqref{eq:ak}, the
  Chebyshev--Wheeler recurrence \eqref{eq:CWrecurrence}, and the orthogonal
  polynomial recurrence \eqref{eq:OPrecurrence} --- are executed over $\C
  [u]$. The result is a rational approximant $y_M (t, u)$ whose numerator and
  denominator coefficients are elements of $\C [u]$, evaluated at a specific
  $u$ only at the last step. This is the operative computational virtue: the
  symbolic pass is performed once at cost $O (M^2)$ polynomial operations in
  $\C [u]$; subsequent evaluations over a Fourier grid in $u$ reduce to
  polynomial substitution.
\end{remark}

\section{Closure Under Integration: The Rough Heston
Application}\label{sec:roughHeston}

\subsection{Rough Heston setup}

In the rough Heston model {\cite{ElEuchRosenbaum2019}}, \eqref{eq:fRiccati} is
instantiated with
\begin{equation}
  P (u) = \tfrac{1}{2} (- u^2 - iu), \quad Q (u) = iu \rho \nu - \lambda,
  \quad R (u) = \tfrac{1}{2} \nu^2, \quad \mu = H + \tfrac{1}{2} \eqqcolon
  \alpha \in (\tfrac{1}{2}, 1), \label{eq:roughHparams}
\end{equation}
where $H \in (0, \tfrac{1}{2})$ is the Hurst exponent, $\lambda > 0$ is mean
reversion, $\nu > 0$ is volatility of volatility, and $\rho \in (- 1, 1)$ is
correlation. All three of $P, Q, R$ lie in $\C [u]$, so 
Theorem~\ref{thm:FRMP} applies directly. Writing $\Lk [x^k] = a (k, u)$ for
the resulting moment functional, the characteristic function is $\varphi_T (u)
= \exp (\kappa (T, u))$ with $\kappa$ as in \eqref{eq:phiT}.

\subsection{Termwise integration}

\begin{lemma}
  [Time-integral M{\"u}ntz coefficients]\label{lem:timeintegral}With $h (t, u)
  = \sum_{k \ge 0} a (k, u)  \hspace{0.17em} t^{(k + 1) \alpha}$,
  \begin{equation}
    \theta \int_0^T h (s, u)  \hspace{0.17em} ds = \theta \sum_{k =
    0}^{\infty} \frac{a (k, u)}{(k + 1) \alpha + 1}  \hspace{0.17em} T^{(k +
    1) \alpha + 1} . \label{eq:timeintegral}
  \end{equation}
\end{lemma}

\begin{proof}
  Integrate termwise: $\int_0^T s^{(k + 1) \alpha}  \hspace{0.17em} ds = T^{(k
  + 1) \alpha + 1} / ((k + 1) \alpha + 1)$. Termwise integration is justified
  by the absolute convergence of \eqref{eq:Muntzseries} on $[0, T_0]$
  established in  Lemma~\ref{lem:Muntz}.
\end{proof}

\begin{lemma}
  [Fractional-integral M{\"u}ntz coefficients]\label{lem:fracintegral}With
  $\alpha \in (\tfrac{1}{2}, 1)$,
  \begin{equation}
    V_0  \hspace{0.17em} I^{1 - \alpha} h (T, u) = V_0  \sum_{k = 0}^{\infty}
    a (k, u) \hspace{0.17em} \frac{\Gamma ((k + 1) \alpha + 1)}{\Gamma (k
    \alpha + 2)}  \hspace{0.17em} T^{k \alpha + 1} . \label{eq:fracintegral}
  \end{equation}
\end{lemma}

\begin{proof}
  Apply the power rule \eqref{eq:powerrule} with $\mu$ replaced by $1 -
  \alpha$ and $s = (k + 1) \alpha$:
  \[ I^{1 - \alpha}  \hspace{0.17em} t^{(k + 1) \alpha} = \frac{\Gamma ((k +
     1) \alpha + 1)}{\Gamma ((k + 1) \alpha + (1 - \alpha) + 1)} 
     \hspace{0.17em} t^{(k + 1) \alpha + 1 - \alpha} = \frac{\Gamma ((k + 1)
     \alpha + 1)}{\Gamma (k \alpha + 2)}  \hspace{0.17em} t^{k \alpha + 1} .
  \]
  Summing over $k$ and multiplying by $V_0$ gives \eqref{eq:fracintegral},
  with termwise application justified as in  Lemma~\ref{lem:timeintegral}.
\end{proof}

\subsection{Reindexing onto the common lattice}

The time-integral sum \eqref{eq:timeintegral} sits on the lattice $\{T^{(k +
1) \alpha + 1} \}_{k \ge 0}$. The fractional-integral sum
\eqref{eq:fracintegral} sits on the shifted lattice $\{T^{k \alpha + 1} \}_{k
\ge 0}$. Setting $k \mapsto k - 1$ in \eqref{eq:fracintegral} aligns it to
$\{T^{(k - 1) \alpha + 1 + \alpha} \} = \{T^{k \alpha + 1} \} = \{T^{(k + 1)
\alpha + 1} |_{k \mapsto k - 1} \}$, i.e., the fractional-integral
contribution at lattice position $k$ involves $a (k - 1, u)$ (with $a (- 1, u)
\assign 0$).

\subsection{The $d (k, u)$ coefficients}

\begin{corollary}
  [$d (k, u)$ coefficient formula]\label{cor:dk}The cumulant generating
  function satisfies
  \begin{equation}
    \kappa (T, u) = \sum_{k = 0}^{\infty} d (k, u)  \hspace{0.17em} T^{(k + 1)
    \alpha + 1}, \label{eq:kappaseries}
  \end{equation}
  where
  \begin{equation}
    \fbox{\concat{d*}{(k,u)}{=<theta>*}{\space{0.17em}}{\frac{\concat{a}{(k,u)}}{\concat{(k+1)}{*<alpha>+1}}}{+V}{\rsub{0}}{*}{\space{0.17em}}{\frac{\concat{<Gamma>*}{(k*<alpha>+1)}}{\concat{<Gamma>*}{((k-1)*<alpha>+2)}}}{*}{\space{0.17em}}{a}{(k-1,u)}{,}{\space{2em}}{k<ge>0,}}
    \label{eq:dk}
  \end{equation}
  with $a (- 1, u) \assign 0$. In particular, $d (k, u) \in \C [u]$ and $d (k,
  u)$ is affinely linear in the pair $(a (k - 1, u), \hspace{0.17em} a (k,
  u))$ with $z$-independent scalar weights depending only on $(\theta, V_0,
  \alpha, k)$.
\end{corollary}

\begin{proof}
  Add \eqref{eq:timeintegral} and \eqref{eq:fracintegral}. Reindex the
  fractional-integral sum $k \mapsto k - 1$: the term formerly at index $k$ in
  \eqref{eq:fracintegral} carries monomial $T^{k \alpha + 1} = T^{(k + 1)
  \alpha + 1 - \alpha}$. Setting $k_{\mathrm{new}} = k + 1$ so the exponent
  becomes $T^{(k_{\mathrm{new}}) \alpha + 1} = T^{((k_{\mathrm{new}} - 1) + 1)
  \alpha + 1}$, the two sums combine at each lattice index $k \ge 0$ to give
  \[ d (k, u) = \theta \hspace{0.17em} \frac{a (k, u)}{(k + 1) \alpha + 1} +
     V_0  \hspace{0.17em} \frac{\Gamma (k \alpha + 1)}{\Gamma ((k - 1) \alpha
     + 2)}  \hspace{0.17em} a (k - 1, u), \]
  where the gamma-ratio arises from \eqref{eq:fracintegral} evaluated at $k -
  1$: $\Gamma ((k - 1 + 1) \alpha + 1) / \Gamma ((k - 1) \alpha + 2) = \Gamma
  (k \alpha + 1) / \Gamma ((k - 1) \alpha + 2)$. Since $a (k, u), a (k - 1, u)
  \in \C [u]$ and the scalar weights are independent of $u$, we have $d (k, u)
  \in \C [u]$. Affine linearity in $(a (k - 1, u), a (k, u))$ is immediate
  from the formula.
\end{proof}

\section{Closure Theorem and the $\kappa$ Approximant}\label{sec:closure}

\begin{theorem}
  [Closure of the M{\"u}ntz--Pad{\'e} apparatus under
  integration]\label{thm:closure}Define the moment functional $\Lkap : \C [x]
  \to \C [u]$ by $\Lkap [x^k] \assign d (k, u)$, where $d (k, u)$ is given by
  \eqref{eq:dk}. Under the quasi-definiteness of $\Lk$, $\Lkap$ is also
  quasi-definite. Let $P_{\kappa} (n, x, u)$ be the orthogonal polynomials of
  $\Lkap$, generated by the Chebyshev--Wheeler recurrence on $S_{\kappa} (n,
  z)$ with $S_{\kappa} (0, z) = \sum_{k \ge 0} d (k, u) z^k$. Define
  \begin{equation}
    \kappa_M (T, u) \assign T \cdot \frac{P_{M, \kappa} (T^{\alpha},
    u)}{T^{\alpha M}  \hspace{0.17em} P_{\kappa} (M, T^{- \alpha}, u)} .
    \label{eq:kappaM}
  \end{equation}
  Then
  \begin{equation}
    \kappa (T, u) = \lim_{M \to \infty} \kappa_M (T, u),
    \label{eq:kappaconvergence}
  \end{equation}
  in the same sense of capacity convergence as  Theorem~\ref{thm:FRMP}.
  Consequently
  \begin{equation}
    \varphi_T (u) = \lim_{M \to \infty} \exp \hspace{-0.17em} (\kappa_M (T,
    u)) . \label{eq:phiconvergence}
  \end{equation}
\end{theorem}

\begin{proof}
  By  Corollary~\ref{cor:dk}, $\kappa (T, u) = T \cdot \sum_{k \ge 0} d (k, u)
  \hspace{0.17em} T^{(k + 1) \alpha}$. Setting $z = T^{\alpha}$, this is $T
  \cdot \tilde{g} (z, u)$ where $\tilde{g} (z, u) = \sum_{k \ge 0} d (k, u) 
  \hspace{0.17em} z^k$ is the generating series of $\Lkap$. Since $d (k, u)$
  is affinely linear in $(a (k - 1, u), a (k, u))$ and the latter satisfy the
  contraction-type bound from  Lemma~\ref{lem:Muntz}, the sequence $\{d (k,
  u)\}$ has the same exponential growth rate as $\{a (k, u)\}$, so $\tilde{g}
  (z, u)$ has positive radius of convergence. Since $d (k, u)$ is affinely
  linear in $(a (k - 1, u), a (k, u))$, the sequence $\{d (k, u)\}$ satisfies
  the same Volterra-iteration structure as $\{a (k, u)\}$, and its generating
  series $\tilde{g}$ likewise satisfies a quadratic functional equation over
  entire functions of $z$ (with the same entire coefficient functions
  $\tilde{P}, \tilde{Q}, \tilde{R}$ of  Lemma~\ref{lem:meromorphic}, linearly
  reweighted by $(\theta, V_0, \alpha)$). Hence $\tilde{g} (z, u)$ extends
  meromorphically to $\C$ by  Lemma~\ref{lem:meromorphic}. Quasi-definiteness
  of $\Lkap$ follows from the affine-linear coupling: the Hankel determinants
  of $\{d (k, u)\}$ are polynomial expressions in those of $\{a (k, u)\}$ and
  hence nonzero generically in $u$. Applying  Theorem~\ref{thm:FRMP} to
  $\tilde{g} (z, u)$ gives \eqref{eq:kappaconvergence};
  \eqref{eq:phiconvergence} follows by continuity of the exponential.
\end{proof}

\begin{corollary}
  [Computational summary for rough Heston]\label{cor:algorithm}Given model
  parameters $(H, \lambda, \nu, \rho, \theta, V_0)$ and truncation order $M$:
  \begin{enumerate}
    \item Compute $a (0, u), \ldots, a (M, u) \in \C [u]$ via
    \eqref{eq:a0}--\eqref{eq:ak} with \eqref{eq:roughHparams}. Cost: $O (M^2)$
    polynomial multiplications in $\C [u]$.
    
    \item Compute $d (0, u), \ldots, d (M, u) \in \C [u]$ via \eqref{eq:dk}.
    Cost: $O (M)$ scalar-polynomial multiplications.
    
    \item Run the Chebyshev--Wheeler recurrence \eqref{eq:CWrecurrence} on
    $S_{\kappa} (0, z) = \sum_{k = 0}^M d (k, u) z^k$ to obtain
    $\alpha_{\kappa} (n), \beta_{\kappa} (n) \in \C [u]$ for $0 \le n < M$.
    Cost: $O (M^2)$ operations in $\C [u]$.
    
    \item Evaluate $\kappa_M (T, u)$ via \eqref{eq:kappaM} at any $(T, u)$ by
    polynomial evaluation. Cost: $O (M)$ per evaluation point.
  \end{enumerate}
  Steps (i)--(iii) are performed {\tmem{once}}, symbolically. Step (iv) is the
  only per-$(T, u)$ cost, reducing to $O (M)$ arithmetic operations regardless
  of the size of the Fourier grid.
\end{corollary}

\section{Structural Summary}\label{sec:structure}

The following table summarises the closure properties established above.

\begin{center}
  \begin{tabular}{llll}
    \hline
    Operation & Lattice & Coefficient at index $k$ & Coupling\\
    \hline
    $h (T, u)$ & $\{T^{(k + 1) \alpha} \}$ & $a (k, u)$ &
    Theorem~\ref{thm:FRMP}\\
    $\int_0^T h \hspace{0.17em} ds$ & $\{T^{(k + 1) \alpha + 1} \}$ &
    $\dfrac{a (k, u)}{(k + 1) \alpha + 1}$ & $a (k)$ only\\
    $I^{1 - \alpha} h$ & $\{T^{k \alpha + 1} \}$ & $\dfrac{\Gamma ((k + 1)
    \alpha + 1)}{\Gamma (k \alpha + 2)} a (k, u)$ & $a (k - 1)$ after
    reindex\\
    $\kappa (T, u)$ & $\{T^{(k + 1) \alpha + 1} \}$ & $d (k, u)$ (eq.
    \eqref{eq:dk}) & affine in $(a (k - 1), a (k))$\\
    \hline
  \end{tabular}
\end{center}

\begin{thebibliography}{99}
  {\bibitem{ElEuchRosenbaum2019}}El~Euch, O. and Rosenbaum, M. {\newblock}The
  characteristic function of rough Heston models.
  {\newblock}\tmtextit{Mathematical Finance}, 29(1):3--38, 2019.
  
  {\bibitem{GatheralRadoicic2019}}Gatheral, J. and Radi{\v c}oi{\'c}, R.
  {\newblock}Rational approximation of the rough Heston solution.
  {\newblock}\tmtextit{Quantitative Finance}, 19(3):495--502, 2019.
  
  {\bibitem{AbiJaberElEuch2019}}Abi~Jaber, E. and El~Euch, O.
  {\newblock}Multifactor approximation of rough volatility models.
  {\newblock}\tmtextit{SIAM Journal on Financial Mathematics}, 10(2):309--349,
  2019.
  
  {\bibitem{Diethelm2004}}Diethelm, K., Ford, N.{\hspace{0.17em}}J., and
  Freed, A.{\hspace{0.17em}}D. {\newblock}Detailed error analysis for a
  fractional Adams method. {\newblock}\tmtextit{Numerical Algorithms},
  36(1):31--52, 2004.
  
  {\bibitem{Gragg1974}}Gragg, W.{\hspace{0.17em}}B. {\newblock}The Pad{\'e}
  table and its relation to certain algorithms of numerical analysis.
  {\newblock}\tmtextit{SIAM Review}, 14(1):1--62, 1972.
  
  {\bibitem{Stahl1997}}Stahl, H. {\newblock}The convergence of Pad{\'e}
  approximants to functions with branch points. {\newblock}\tmtextit{Journal
  of Approximation Theory}, 91(2):139--204, 1997.
  
  {\bibitem{Chihara1978}}Chihara, T.{\hspace{0.17em}}S.
  {\newblock}\tmtextit{An Introduction to Orthogonal Polynomials}.
  {\newblock}Gordon and Breach, New York, 1978.
  
  {\bibitem{MuntzLegendre2015}}Esmaeili, H. and Shamsi, M. {\newblock}A
  pseudo-spectral scheme for the approximate solution of a family of
  fractional differential equations. {\newblock}\tmtextit{Applied Mathematical
  Modelling}, 39(16):5148--5163, 2015.
  
  {\bibitem{WheelerGautschi1996}}Gautschi, W. {\newblock}\tmtextit{Orthogonal
  Polynomials: Computation and Approximation}. {\newblock}Oxford University
  Press, Oxford, 2004.
\end{thebibliography}

\end{document}
